% Compiles on its own (pulls ../main.tex's preamble via subfiles) or as part of the book.
\documentclass[../main.tex]{subfiles}
\begin{document}

\chapter{One Image Through the Pipeline}

Theory is easier to trust when you can watch it happen to something real. This chapter follows a
single KAIST frame --- \texttt{set06/V000/I01259}, a daytime tree-lined avenue with one pedestrian
--- through every step built so far, and looks at what each step does to it. It is a living
chapter: a panel is added as the pipeline grows (neck, head, fusion, and finally the same frame
after training).

\begin{figure}[htbp]
  \centering
  \includegraphics[width=\linewidth]{pipeline_trace_phase1}
  \caption{One real KAIST frame through PRISM so far. \textbf{Row 1 (load):} the RGB and thermal
  pair with the ground-truth box, at identical pixel coordinates in both. \textbf{Row 2
  (augment):} the same pair after a paired horizontal flip --- both streams and the box move
  together. \textbf{Row 3 (backbone):} the RGB image squeezed into its three feature-map summaries
  \texttt{dark3}/\texttt{dark4}/\texttt{dark5}, shown as activation heatmaps. Regenerate with
  \texttt{python scripts/trace\_image.py}.}
  \label{fig:trace}
\end{figure}

\section{Row 1 --- Load}

The frame arrives as a \emph{pair}: the visible RGB image and the thermal (LWIR) image, both
looking at the same scene. The green box marks the one scored pedestrian, read from the sanitized
annotation. The box sits at the \emph{same pixel coordinates in both images} --- and that shared
alignment is precisely what fusion will later depend on. Behind the picture, the loader has
converted the box from $x,y,w,h$ to $x_1,y_1,x_2,y_2$, marked \texttt{person} as scored and the
other classes as ignore regions, and returned two $640\times512$ tensors.

\section{Row 2 --- Augment}

The same pair after a paired horizontal flip. The pedestrian, the lane markings, everything moves
to the opposite side --- and it moves \emph{identically in both streams}, with the box following
the pixels. This is the augmentation rule from Chapter~2 made visible: geometric operations are
applied the same way to both modalities, so the pair never comes unstuck. The correspondence unit
test is the automatic proof that this property holds; the figure is the human-readable version of
the same guarantee.

\section{Row 3 --- Backbone}

The RGB image squeezed into its three summaries, shown as heatmaps (brighter = stronger
activation). Read them left to right and watch the resolution collapse:
\begin{itemize}
  \item \texttt{dark3} --- $64\times80$, stride 8. Still detailed: road edges, trees, and lane
        lines are visible.
  \item \texttt{dark4} --- $32\times40$, stride 16. Coarser; fine detail is gone, broad structure
        remains.
  \item \texttt{dark5} --- $16\times20$, stride 32. Very coarse --- bright region here, road there.
        It knows roughly \emph{what} is \emph{where}, not precisely where.
\end{itemize}

Two honest points about this row.

First, the backbone is \emph{untrained} --- its weights are random. So the maps highlight edges and
texture, not pedestrians; after training they will fire on people instead. The point of the row
right now is not what the maps detect but the \emph{spatial shrink}.

Second --- and this is the payoff --- that shrink is the misalignment story from Chapter~1, drawn.
An $8$-pixel offset between the RGB and thermal cameras is a full cell at \texttt{dark3}
(stride 8) but only a quarter of a cell at \texttt{dark5} (stride 32). So \emph{where} we fuse
directly sets how tolerant the detector is to the two cameras not lining up --- the entire Phase-4
headline is visible here, in the coarsening from left to right.

These three maps are exactly where fusion attaches. When we go two-stream, a second backbone
produces a second \texttt{dark3}/\texttt{dark4}/\texttt{dark5} from the thermal image, and the
fusion operator merges the two sets before the neck ever sees them.

\end{document}
